\section{COVER-EVO Method}

\subsection{Fixed Construction and Pruning}

COVER-EVO uses a fixed greedy construction procedure in
which only the vertex-scoring function is evolved. Starting
with an empty dominating set $D=\varnothing$ and uncovered
vertices $U=V$, the algorithm repeatedly selects the
unselected vertex $v$ with the highest score $f(v,U)$
among those satisfying $|N[v]\cap U|>0$. The selected
vertex is added to $D$, and its closed neighborhood is
removed from $U$. Ties are resolved by lower vertex cost,
then lower vertex ID. Previously covered vertices remain
eligible if they can cover additional uncovered vertices.

Once all vertices are covered, a deterministic pruning
pass examines selected vertices in descending cost order,
breaking ties by ascending ID. A vertex is removed only
if the remaining set still dominates the graph. This
construction and pruning procedure is shared by all
scoring-based methods.

Feasibility and objective values are independently
verified against the original graph. Operational failures
receive a penalty equal to the cost of selecting all
vertices and remain included in evaluation statistics.

\subsection{Scoring Features and DSL}

For a candidate vertex $v$, define its newly covered
vertices as $A_v=N[v]\cap U$, coverage gain
$k_v=|A_v|$, and mean vertex cost
$\bar c=n^{-1}\sum_{v\in V}c_v$.
The scoring language exposes five features:
\begin{align}
g_v &= \frac{k_v}{|U|},
&q_v &= \frac{c_v}{\bar c}, \notag\\
r_v &= \frac{1}{k_v}\sum_{u\in A_v}
       \frac{1}{|N[u]|},
&o_v &= 1-\frac{k_v}{|N[v]|},
\label{eq:features}\\
t &= \frac{|U|}{n}. &&
\notag
\end{align}

Here, $g_v$ measures normalized coverage gain, $q_v$
normalized selection cost, and $r_v$ the average rarity
of newly covered vertices. The overlap feature $o_v$
measures the fraction of the candidate's neighborhood
already covered, while $t$ represents the remaining
uncovered fraction. Rarity is computed using original
neighborhood sizes rather than dynamically updated
degrees.

\begin{samepage}
Programs are scalar expressions over these features and bounded constants,
using arithmetic, protected division, minimum, maximum and absolute value.
Control flow, assignments, imports and external calls are forbidden.
The repository specifies expression-size limits and numerical safeguards.\par
\end{samepage}

The initial greedy seed is
$f_0(v,U)=g_v/q_v$, which induces the same ranking
as $k_v/c_v$. I additionally consider a handcrafted
rarity-weighted rule:
\begin{equation}
f_{\mathrm{RGR}}(v,U)
    =\frac{g_vr_v}{q_v}.
\label{eq:rarity}
\end{equation}

GR+P denotes standard weighted greedy with pruning,
while RGR+P uses Eq.~\eqref{eq:rarity} with the same
construction and pruning procedure. RGR+P is a rarity-weighted coverage
priority; Section~\ref{sec:rarity_form} derives its equivalence to an
inverse-neighborhood-weighted demand sum per unit cost. Candidate
priorities are evaluated using exact rational arithmetic.
Reverse-delete provides an additional baseline by
starting from all vertices and applying the pruning
procedure directly.

\subsection{LLM-Guided Discovery}

Each search run begins with the greedy seed and evaluates candidate
programs on the training graphs using mean cost ratio to GR+P.
An archive maintains the four best candidates, ranked
by training fitness, AST size, and program hash.
Every proposal attempt counts toward the budget without replacement.

\textbf{IND} generates candidates independently,
without incorporating previous programs or evaluation
results into the generation prompt. It nevertheless
uses the same evaluator, archive, and final selection
procedure as the evolutionary conditions.

\textbf{EVO-P} conditions new proposals on up to two
sampled archive parents, providing their expressions
and rationales without explicit numerical performance
feedback.

\textbf{EVO-F} follows the same parent-conditioned
procedure but additionally provides numerical training
summaries and descriptors of poorly performing cases.
Both evolutionary conditions begin with eight fresh
proposals and subsequently include one fresh proposal
after every four parent-conditioned attempts.

EVO-P still uses training fitness for archive selection. EVO-F versus
EVO-P therefore compares a prompt-feedback policy with a parent-only
policy, both using fitness-based selection. It adds numerical summaries
and case descriptors, while subsequent archives and parent exposure can
diverge. The contrast concerns complete proposal policies.
EVO-F versus IND compares the combined parent-conditioned and
feedback-guided policy with independent generation.

\subsection{Controls and Winner Selection}

To assess whether simple parameter variation can
produce competitive heuristics, I include RS-PARAM,
a random search over the scoring family
\begin{equation}
f_{a,b,d}(v,U)=
\frac{g_v}{q_v}
(1+ar_v+bo_v+dr_vt).
\label{eq:rs}
\end{equation}

Coefficients $a,b,d$ are independently sampled from
$\{-2.000,-1.999,\ldots,2.000\}$. RS-PARAM uses
the same proposal-attempt budget, training evaluator,
greedy seed, and validation procedure, without LLM
calls. Its parameterized search space is not identical
to the full expression language.

Following discovery, the three highest-ranked training
candidates from each run are evaluated on the validation
graphs. One winner is selected per run and frozen before
held-out testing. Validation results are never supplied
to the proposal model.
